%%%%%%%%%%%%%%%%%%%%%%%%%%%%%%%%%%%%%%%%%%%%%%%%%%%%
\section{Algorithm}

\subsection{Monotone Scan}

For a fixed itemset $X$, let $\mathbf{p} = (p_0, p_1, \ldots)$ be the sorted
list of timestamps at which $X$ occurs.
We scan the left endpoint $l$ from $0$ to the current position,
maintaining a pointer into $\mathbf{p}$.

\textbf{Key property}: as $l$ increases by one, the window $[l, l+W)$
loses $T_{l-1}$ and may gain $T_{l+W-1}$.
Consequently, $\mathrm{freq}(X, l)$ changes by at most $\pm 1$ per step,
and the left endpoint $l$ is monotonically non-decreasing.
This means any confirmed output is never invalidated.

\subsection{Stride-Based Skipping}

Na\"ively incrementing $l$ by 1 each step is wasteful.
We exploit two cases:

\textbf{Dense phase}: if $\mathrm{freq}(X, l) \geq \theta$,
the next non-dense position is at least $l + \mathrm{surplus}$, where
$\mathrm{surplus} = \mathrm{freq}(X, l) - \theta + 1$.
We skip ahead by $\mathrm{surplus}$ steps.

\textbf{Non-dense phase}: if $\mathrm{freq}(X, l) < \theta$,
the next potentially dense position requires
$\theta - \mathrm{freq}(X, l)$ additional occurrences.
We advance $l$ to the position of the $(\theta - \mathrm{deficit})$-th
next occurrence, skipping the deficit gap.

\subsection{Apriori-Style Candidate Pruning}

For itemset $X$ of size $k$, if any $(k{-}1)$-subset $Y \subset X$
has no dense interval intersecting $[l, l+W)$, then $X$ cannot be
dense in $[l, l+W)$ either (by the anti-monotonicity of support).
We precompute, for each $(k{-}1)$-subset, the set of $l$-values that
are dense, and restrict candidate scan ranges for $X$ to the
intersection of its subsets' dense ranges.

\subsection{n-Step Streaming}

In the target deployment, transactions arrive continuously from a
sensor stream or autonomous agent log and are appended to a database.
The system waits until $n$ new transactions have accumulated since the
last scan, then triggers a mining step.
The parameter $n$ directly controls output latency: smaller $n$ yields
more responsive output at higher computational cost; larger $n$ batches
more events per scan, reducing cost.
Setting $n = T$ (full stream length) recovers batch processing as a
special case.

Each scan step proceeds as follows:

\begin{enumerate}
  \item When $n$ new transactions have arrived (current position $s$),
        retrieve the buffer $[s - W + 1,\; s + n)$, retaining the
        $W{-}1$ most recent transactions from the previous step to
        prevent boundary loss.
  \item Run the monotone scan with stride-based skipping and
        Apriori-style pruning over this buffer.
  \item Emit all intervals $(X, [l_s, l_e])$ with $l_e + W \leq s$
        not yet output: these are \textbf{confirmed} and will not
        change as future transactions arrive.
  \item Advance to the next accumulation of $n$ transactions and repeat.
\end{enumerate}
